\odiseachapter{esqueria}{Esqueria}\label{cap:esqueria}

Cuando por fin Odiseo se decide a salir del bosque y dirigirse a palacio, Atenea se deja de remilgos y vuelve a la carga ---esta mujer no se está nunca quieta---. Primero se hace la encontradiza, fingiéndose una niña, y luego guía al héroe hasta la residencia de Alcínoo:

\begin{versocita}
A la que dijo ella así, por delante fue Palas Atena\\
rauda guiándolo, y él tras la diosa pisando sus huellas.\\
Los navicélebres feacios tampoco se daban ni cuenta\\
de que anduviera por la ciudad entre ellos; ni Atena\\
lo permitía, la diosa zahorí, trenzabella, que en niebla\\
de maravilla envolvía al ser de su benevolencia.\\
Y, boquiabierto, Odiseo los puertos, las naos niveleñas,\\
las tantas plazas de próceres, y las murallas contempla,\\
largas, altivas, con picas ---asombro de ver--- de acroteras.\\
Cuando a la muy alta casa del alto señor ellos llegan,\\
así la diosa, ojiglauca Atenea, le abría conversa:
\end{versocita}

Observa, atenta lectora, exacto lector, cómo Homero nos dibuja, en unas pocas líneas, la próspera y marina tierra de los \emph{navicélebres} feacios. Los puertos, las plazas, las naos niveleñas, las largas y altivas murallas. El exiliado, boquiabierto, comprende que ha llegado a un lugar civilizado, quizás el lugar más civilizado que ha visitado nunca. Si Nausícaa puede simbolizar el amor imposible, Esqueria, la tierra de los feacios, representa el hogar que supera a Ítaca en todo, menos en la nostalgia.

Haz algo más. Lee en voz alta ---a ser posible gritando y con grandes aspavientos--- los versos. Escucha el sonido de tu voz, como música. Estamos en la tierra de los feacios y el Poeta se hace eco de su belleza con versos aún más bellos, que Juanma Tobal capta con su oído certero. ¡Oh diosa zahorí, trenzabella, que en niebla
de maravilla envolvía al ser de su benevolencia!

A la puerta de palacio, Atenea, todavía disfrazada, le larga un tremendo discurso a Odiseo informándole de toda la genealogía de Alcínoo y su esposa Arete antes de marcharse. Nuestro héroe entra en palacio y se maravilla de todas las riquezas que encuentra. Homero no escatima detalles:

\begin{versocita}
Que era el fulgor de la luna o del sol lo que allí relumbraba\\
por esas salas altivas de Alcínoo buenasentrañas.\\
Y es que paredes de bronce corrían a banda y a banda\\
desde la entrada hasta el fondo, en un friso de azul rematadas;\\
de oro el portón que cerraba por dentro la sólida casa,\\
y las sus jambas de plata, al suelo bronceño clavadas,\\
de plata encima el dintel, y oridorada la aldaba.
\end{versocita}

La descripción de la fabulosa morada continúa por un buen rato y uno tiene la sospecha de que la casa de Ulises en Ítaca e incluso el palacio de Menelao en Esparta son chozas cabreras comparadas con todo el esplendor que al náufrago se le ofrece. El Poeta, que a contable no le gana nadie, no solo da una descripción minuciosa de ornamentos, oros, platas y bronces, sino que también hace recuento de esclavos, sirvientes, telas, perales, granados, manzanos, higueras, olivos y viñedos que dan exquisito vino. Esqueria, en una palabra, es Jauja.

El caso es que el héroe, siempre envuelto en la niebla de Atenea, se las compone para llegar al trono y allí abraza las rodillas de Arete y le dice:

\begin{versocita}
«Hija del pardelosdioses Rexénor, Arete, de hinojos\\
a tu señor, tus rodillas, tras tanto pasado, me postro\\
y a los aquí convidados. Vivir os dé el cielo gozosos,\\
y que entregar en legado a los hijos podáis un día todos\\
de vuestra casa el caudal y el honor que os dio el pueblo en decoro.\\
Aparejadme una escolta que a casa me lleve, os lo imploro,\\
pronto, que ha mucho que sufro, y estoy sin los míos y solo».
\end{versocita}

Fíjate bien, compasiva lectora, perdonatodo lector. Este hombre que ha saqueado Troya se ve reducido a la condición de mendigo, se ve obligado a implorar donde antes se habría limitado a arrebatar. En Esqueria descubre ---¿o finge?--- la humildad, aprende a suplicar, ve, quizás por vez primera, los toros de puntiagudos cuernos desde el otro lado de la barrera. Esqueria no solo es la quintaesencia de la hospitalidad. Es el lugar donde Odiseo se vuelve algo más humano.

Y Nolan la elimina de su historia.

¿Por qué? El tema central de la \emph{Nodisea} puede resumirse así: la ley de Zeus obliga a los hombres a la \gr{ξενία} (\emph{xenía}), esto es, a ser hospitalarios ---no por virtud, por cierto, sino por precaución, ya que el forastero al que se recibe bien podría ser un dios disfrazado---. Ulises y los aqueos toman Troya con un engaño, saltándose, según afirma el osado director, esa ley y precipitando la caída de la civilización y el final de la Edad del Bronce. Es una tesis un poco extrema, sobre todo porque aqueos y troyanos están en guerra y la treta del caballo no involucra una petición de hospitalidad sino un abuso del sentimiento religioso ---o más plausiblemente del orgullo de los de Ilión que se sienten vencedores---. Pero el caso es que nuestro fílmico caballero se aferra a ella como a un clavo ardiendo. Quizás por eso suprime de la película uno de los episodios más importantes de la \emph{Odisea}, ya que Nausícaa y los feacios suponen un contraejemplo a su pesimismo. Y en mi opinión eso es un pecado capital, que puede suponerle a sir Christopher varios siglos de purgatorio. No solo elimina uno de los personajes más adorables ---quizás el más adorable, en todos los sentidos--- de la épica de Homero, sino que borra las pruebas que no le convienen a su narrativa.

Después de abrazar las rodillas de la reina, Ulises se dirige al hogar, donde se sienta «sobre las cenizas, cerca del fuego».
El viejo Equeneo le dice entonces al rey:

\begin{versocita}
«Alcínoo, no es para ti ni nobleza ni bien que ese huésped\\
ante el fogón, por sobre la ceniza, en el suelo se siente;\\
y, porque esperan tu voz, todos estos aquí se contienen.\\
Haz levantar a ese hombre y asiento en silla tú ofrécele\\
claveteada de plata, y manda también que nos mezclen\\
más vino aún los heraldos, que ahora al que juega a que truene\\
le arramaremos, a Zeus, que acompaña al que en buen ruego viene;\\
y al extranjero comida dé el ama de lo que almacene».
\end{versocita}

El buen rey no se hace de rogar, recoge al mendigo, lo lleva a su mesa, lo alimenta y lo agasaja. Luego despide a todo el mundo, no sin avisar de que al día siguiente planearán el regreso del huésped. Y aquí aparecen las célebres líneas que justifican formalmente la \gr{ξενία}, esa hospitalidad debida a los extraños:

\begin{versocita}
Y si resulta que él es un sin muerte que el cielo aquí manda,\\
señal será de que el cielo alguna otra cosa se trama.\\
Que ante nosotros los dioses han hecho visible su estampa\\
siempre que se hizo en su honor ciemboyada gloriosa, y es fama\\
que a nuestro lado han comido sentándose en nuestras bancadas.
\end{versocita}

Odiseo, por su parte, responde:

\begin{versocita}
«Sea otro, Alcínoo, el cuidado en que des en pensar; no me veo\\
yo como aquellos sin muerte que pueblan la campa del cielo\\
ni en porte ni en condición, y sí un hombre de los morideros.\\
A ese cualquiera que veis soportar de más bregas el peso,\\
a ese hombre sí soy igual en cosa de sufrimiento.\\
\stanzabreak
Y hasta podría yo hacer de desdichas más grande recuento,\\
cuantas ---y bien que lo sé--- por querer de los dioses yo llevo.\\
Mas permitid que ahora cene, dolido y todo que vengo;\\
cosa más perra no hay que ese impúdico estómago nuestro\\
que nos reclama y apremia a tenerlo a la fuerza en recuerdo\\
por agobiados que estemos y penas haya en el pecho;\\
sí, que mil penas yo llevo en el pecho, y él anda pidiendo\\
siempre que coma y que beba, y que ya no me acuerde de eso,\\
de lo que tanto he pasado, y me llama a dejarlo bien lleno.\\
\stanzabreak
Pero vosotros, del alba a las claras, haced ---os lo ruego---\\
por prepararle a este pobre infeliz a su patria el regreso,\\
con sus pesares a cuestas; que no se me marche mi tiempo\\
sin ver mi hacienda otra vez, la casa techalta y los siervos».
\end{versocita}

Sin desperdicio la reflexión que se marca sobre las exigencias del estómago. Tampoco lo tiene oírle llamarse «pobre infeliz».

Cuando se marchan los demás y Odiseo queda a solas con Alcínoo y Arete, la reina, muy aguda, se fija en las ropas que lleva y deduce, sin decirlo, que se las ha prestado Nausícaa. Dice entonces:

\begin{versocita}
«Bien, lo primero, extranjero, te pido yo misma respuesta:\\
¿quién y de qué gentes tú? ¿Quién te ha dado esas ropas que llevas?\\
¿No dices tú que del mar, errabundo infeliz, aquí llegas?».
\end{versocita}

Ulises cuenta cómo llegó a Ogigia, la morada de Calipso, después de que Zeus hundiera su barco, cómo pasó allí siete años y cómo el viaje en balsa terminó en naufragio por intervención de Poseidón. Luego dice la verdad en lo que respecta a Nausícaa:

\begin{versocita}
Se hundía el sol cuando, al fin, me soltó el dulce sueño. Mis ojos\\
vieron entonces allí a unas esclavas jugando en el soto,\\
las de tu hija, y tu hija con ellas, diosa en su coro.\\
Le supliqué; y errada no fue en su talante piadoso,\\
ese que nunca esperaras que un joven que encuentras de pronto\\
te regalara: los jóvenes siempre son faltos de aplomo.\\
Y es que ella a mí me dio pan, me dio vino ardinegro ---y a colmo---,\\
y me bañó ella en el río, y me dio el atavío que porto.\\
Hondodolido yo estoy, mas te he dicho verdad, verdad solo.
\end{versocita}

A lo que responde Alcínoo:

\begin{versocita}
«En algo, huésped, sí erró en aquel su talante mi hija,\\
sí, que con sus camareras aquí, a nuestra casa, tendría\\
ella que haberte traído, pues que a ella el ruego le hacías».
\end{versocita}

Y Ulises:

\begin{versocita}
«No me la riñas por eso, señor, a esa justa doncella,\\
que a mí ella sí me invitó a seguirla con sus asistentas,\\
pero le dije que no, por temor y también por vergüenza;\\
vete a saber si, en viéndome tú, mala mueca me hicieras,\\
que muy celosa es la raza de quienes pisamos la tierra».
\end{versocita}

«Que muy celosa es la raza de quienes pisamos la tierra». Hay que decir que Odiseo ---es decir, Homero--- entiende muy bien la condición humana.

Y en este punto Alcínoo nos sale por peteneras. Recuerda, lectora de agudos ojos, lector de penetrante mirada, que todo lo que sabe el rey del mendigo que tiene delante es que se trata de un pobre hombre, un náufrago. Pero, o bien porque la referencia a Calipso le ha dado una firme pista ---pocos mortales se pasan siete años compartiendo casa y lecho con una ninfa---, o bien porque es un hombre acostumbrado a juzgar más allá de las apariencias, el rey ha decidido al vuelo que el extranjero es un buen partido y así se lo hace saber:

\begin{versocita}
A lo que ya contestaba Alcínoo de nuevo diciendo:\\
No llevo yo un corazón, extranjero, tal en el pecho\\
que se enfurruñe tan sin razón: cada cosa en su asiento.\\
Así Zeus Padre y Atena y Apolo tuvieran dispuesto\\
que siendo tú tal cual eres, y nuestros pensares gemelos,\\
a esta hija mía tomaras y fueras llamado mi yerno
\end{versocita}

Pero acto seguido le dice también que, si prefiere volver a su casa, hará todo lo que esté en su mano para devolverlo a su tierra.

\begin{versocita}
[...] casa y bienes pondríalos yo si, queriendo,\\
te nos quedaras; que aquí, sin querer, no te va a poner freno\\
ninguno de los feacios, que no iba a gustar eso a Zeus.\\
\stanzabreak
Tiempo sí pongo a tu guarda y escolta, y querrás bien saberlo,\\
será mañana: rendido tú entonces allí por el sueño,\\
te dormirás, y a tu patria y tu casa remando en mar quieto\\
te llevarán, o allá donde digas que tienes apego,\\
incluso si eso supone llegar hasta Eubea y más lejos,\\
la tierra más retirada nos dicen los que un día la vieron\\
de nuestros hombres, la vez en que a aquel Radamantis el güero\\
diéronle escolta para visitarse con Ticio Terrero.\\
\stanzabreak
Sí, que allí fueron y sin fatigarse en un día cumplieron,\\
y regresaron en ese día mismo a casa de nuevo.\\
Por ti sabrás tú también cuán magníficos son mis veleros\\
y esos muchachos que el vómito arrancan al mar con su remo.
\end{versocita}

¡Magia de nuevo! Hasta aquí, podíamos conformarnos con la idea de que los feacios son, simplemente, buenos marinos. Pero Alcínoo le asegura a Odiseo que sus barcos pueden llevarle a Eubea ---como quien dice el fin del mundo--- en un solo día \emph{y regresar en esa misma jornada}.

Hay una correspondencia encantadora con la visita de Telémaco a Esparta. Terminada la conversación, Arete en persona se ocupa de acomodar a Odiseo ---por supuesto bajo el soportal, el mismo lugar donde Helena había acomodado a Telémaco--- y, una vez que ha cumplido con su deber, Alcínoo y Arete hacen exactamente lo mismo que Helena y Menelao:

\begin{versocita}
Alcínoo, en cambio, en lo hondo de la alta mansión se acostaba;\\
con él su esposa de ley, aliñándole el lecho, se encama.
\end{versocita}
